\centerline{\bf \Huge Великорусский зажим между квадратами}

\begin{flushright}
        \scriptsize{}
\end{flushright}

\textbf{Утверждение.} Даны натуральные числа $a$ и $b$. Если выполнено неравенство $a^2<b<(a+1)^2$, то число $b$ не является точным квадратом.

\problem{1}
Может ли точный квадрат быть равен произведению двух последовательных натуральных чисел?

\problem{2}
Найдите все такие натуральные $n$, для которых $n^2+4 n-1$ является точным квадратом.

\problem{3}
Найдите все пары натуральных чисел $m$ и $n$, при которых оба числа $m^2+2 n$ и $n^2+2 m$ являются точными квадратами.

\problem{4}
Пусть $a_1, a_2, a_3, \ldots-$ бесконечная возрастающая арифметическая прогрессия натуральных чисел. Докажите, что не могут все члены этой прогрессии оказаться точными квадратами.

\problem{5}
Найдите все пары натуральных ( $a, b$ ), для которых числа $a^2+4 b+4$ и $b^2+4 a+4$ являются точными квадратами.

\problem{6}
Паша выбрал целые числа $x$ и $y$. Виктор заметил, что какие бы натуральные числа $k$ и $m$ не взять, число $x k^2+y m^2$ оказывается точным квадратом. Докажите, что $x$ либо $y$ равен нулю.